\answer{ex:01:1}{(а)~$8\cdot10^{10}$ \nt{ГЛУТ} --- $10$ населений Земли; широковещательных $\approx1.9\cdot10^6$ --- город размером с Вену. (б)~$\approx8\cdot10^{14}$ окончаний, из них широковещательных $\approx3\cdot10^{11}$, доля $\approx4\cdot10^{-4}$ --- в $\approx16$ раз больше доли этих нейронов ($2.2\cdot10^{-5}$).}
\answer{ex:01:2}{(а)~$\tau_c\ln(c_0/c_*)\approx4.6\ms$. (б)~Линия свободна при $T\ge\tau_c\ln101$: до $\approx22$ выбросов в секунду вместо $\approx220$.}
\answer{ex:01:3}{$V(s_k)=\gamma^{K-1-k}r$; $\delta_{\text{лампа}}=\gamma^Kr=re^{-T/\tau_\gamma}\approx0.60\,r$ при любом $\Delta$, $\tau_\gamma\approx1.95\,\text{с}$.}
\answer{ex:01:4}{$n^*=93$, $47$, $25$, $12$: $n^*\approx5/\alpha$ при малых $\alpha$.}
\answer{ex:01:5}{\nt{ГЛУТ}: $1\to10\to10^4\to10^7$ --- $\approx10^7$ нейронов, $4$ звена, $4\ms$. \nt{ДОФА}: $1\to10\to3.2\cdot10^5$ --- $\approx3\cdot10^5$ нейронов, $3$ звена, в $30$ раз меньше; тот же порядок, что число \nt{ДОФА}-нейронов.}
\answer{ex:01:6}{$\lambda=f_EJ_E-f_IJ_I$, откуда $J_I=37.5$; отношение токов $f_IJ_I/f_EJ_E=0.94$; лавина ($\lambda\ge1$) при ослаблении торможения всего на $6.7\,\%$.}
\answer{ex:01:7}{$n\ge57$ предъявлений на каждую задержку в каждом состоянии. Тот же спад дала бы, например, неточность внутренней оценки времени, растущая с $T$.}
